%==============================================================================
% BANCO DE PREGUNTAS Y RESPUESTAS PARA LA SUSTENTACIÓN
% Tratamiento Numérico de la KHI bajo el Marco de la RRMHD
% Compilación: pdflatex banco_preguntas.tex (×2)
%==============================================================================
\documentclass[11pt,letterpaper]{article}

\usepackage[utf8]{inputenc}
\usepackage[T1]{fontenc}
\usepackage[spanish,es-tabla]{babel}
\usepackage{lmodern}
\usepackage[margin=2.3cm]{geometry}
\usepackage{xcolor}
\usepackage{booktabs}
\usepackage{enumitem}
\usepackage{amsmath,amssymb}
\usepackage{fancyhdr}
\usepackage{lastpage}
\usepackage[hidelinks]{hyperref}

\definecolor{primary}{HTML}{1B0C41}
\definecolor{accent}{HTML}{E65D2F}
\definecolor{grisnota}{HTML}{666666}

\pagestyle{fancy}
\fancyhf{}
\fancyhead[L]{\footnotesize\itshape Banco de preguntas --- KHI en RRMHD}
\fancyhead[R]{\footnotesize Martinez \& Rodriguez}
\fancyfoot[C]{\footnotesize Página \thepage\ de \pageref{LastPage}}
\renewcommand{\headrulewidth}{0.4pt}

\setlength{\parskip}{3pt}
\setlength{\parindent}{0pt}

% P: pregunta con etiqueta [nivel · apoyo]; R: respuesta
\newcounter{preg}
\newcommand{\Q}[2]{\par\addvspace{11pt}\stepcounter{preg}%
  \noindent{\color{primary}\bfseries P\thepreg.}\ \textbf{#2}\par
  \nopagebreak{\scriptsize\color{grisnota}#1}\par\nopagebreak\vspace{1pt}}
\newcommand{\A}[1]{\noindent\emph{R.}\ #1\par}

\begin{document}

%---------------------------------------------------------- portada
\begin{titlepage}
\centering
\vspace*{1.2cm}
{\color{primary}\rule{\linewidth}{2pt}}\\[0.9cm]
{\large Universidad Distrital Francisco José de Caldas}\\[2pt]
{\small Facultad de Ciencias Matemáticas y Naturales --- Programa Académico de Física}\\[1.6cm]
{\LARGE\bfseries Banco de Preguntas y Respuestas}\\[0.5cm]
{\large Preparación de la Sustentación del Trabajo de Grado}\\[0.9cm]
{\color{primary}\Large\bfseries Tratamiento Numérico de la Inestabilidad de
Kelvin--Helmholtz bajo el Marco de la Magnetohidrodinámica Relativista
Resistiva}\\[1.4cm]
\begin{tabular}{rl}
\textbf{Autores:}  & Bryan Martinez Anzola \;$\cdot$\; Sebastián Rodriguez Garcia\\[4pt]
\textbf{Director:} & Prof.\ Ph.D.\ Sergio Miranda-Aranguren\\[4pt]
\textbf{Jurado:}   & Prof.\ Juan Carlos Giraldo Acuña\\[4pt]
\textbf{Base:}     & monografía (167 pp), presentación (50 láminas) y las\\
                   & 44 referencias citadas\\
\end{tabular}\\[1.5cm]
{\color{primary}\rule{\linewidth}{2pt}}\\[0.8cm]
{\small Bogotá, D.C.\ --- Julio de 2026}
\vfill
\end{titlepage}

\tableofcontents
\newpage

\section*{Cómo usar este banco}
\addcontentsline{toc}{section}{Cómo usar este banco}
Cada pregunta lleva una etiqueta con el \emph{nivel} esperado de exigencia
(básico $=$ verificación de comprensión; medio $=$ dominio del trabajo;
alto $=$ frontera del trabajo o de la literatura) y, cuando aplica, la
\emph{lámina} de la presentación o la lámina de \emph{respaldo} (B1--B13)
que sirve de apoyo visual al responder. Las respuestas están redactadas en
la extensión adecuada para el turno de preguntas: entre treinta segundos y
un minuto de exposición oral. Las categorías siguen el orden de la
sustentación; las dos últimas cubren la dimensión epistemológica ---la de
mayor probabilidad dado el perfil del jurado--- y la bibliografía citada.
Conviene ensayar en voz alta, en particular, todas las preguntas de nivel
básico (fallarlas cuesta más que acertar las de nivel alto) y las marcadas
en las categorías de epistemología y disipación numérica.

%========================= PARTE A: categorías 1–5 ===========================

\section{Fenomenología y contexto astrofísico}

\Q{nivel básico · lámina 3}{Explique la inestabilidad de Kelvin--Helmholtz
sin usar ecuaciones.}
\A{Dos flujos en contacto a velocidades distintas almacenan energía libre en
su capa de cizalla. Si la interfaz se ondula, el flujo se acelera sobre las
crestas y se frena en los valles; por Bernoulli, esa asimetría genera
gradientes de presión que ejercen un torque sobre la interfaz y amplifican
la ondulación. La retroalimentación produce crecimiento exponencial (fase
lineal) y, al saturar, la interfaz se enrolla en vórtices que mezclan ambos
medios (fase no lineal).}

\Q{nivel básico · lámina 3}{¿Dónde opera esta inestabilidad en el universo
relativista y por qué es importante?}
\A{En toda frontera entre un flujo relativista y su entorno: bordes de
chorros de AGN y GRBs, vientos de púlsares (Bucciantini 2006), eyecciones
de objetos compactos, e incluso la magnetopausa terrestre en el régimen
clásico (Hasegawa 2004). Gobierna la mezcla y el frenado del chorro
(\emph{entrainment}), participa en la dicotomía morfológica FR-I/FR-II,
actúa como dínamo local amplificando $B$ (Pimentel 2019) y crea las
láminas de corriente donde la reconexión acelera partículas.}

\Q{nivel medio · lámina 10}{¿Qué diferencia a la KHI relativista de la
clásica?}
\A{Tres cosas. La inercia efectiva es $\rho h W^2$: el factor de Lorentz y
la entalpía aumentan la ``masa'' de la capa y modifican la tasa. La
compresibilidad estabiliza: en el régimen clásico la discontinuidad es
incondicionalmente inestable, mientras que en el relativista rige el
criterio $0<M_{re}<\sqrt2$ con el Mach relativista construido con la
velocidad de restitución apropiada (Königl 1980; Bodo 2004; Chow 2023). Y
la señal se propaga con cota $c$, lo que acopla los sectores
electromagnético e hidrodinámico de manera no galileana.}

\Q{nivel medio · lámina 4}{¿Por qué la reconexión magnética que muestran los
códigos ideales es un problema conceptual?}
\A{Porque en RMHD ideal el teorema de Alfvén prohíbe cambiar la topología
de $\mathbf B$: si un código ideal reconecta, lo hace por el error de
truncamiento de la malla, una ``resistividad numérica'' que depende de la
resolución y no de la física (Komissarov 2007). El resultado no es
predictivo: cambiar la malla cambia la tasa de reconexión. La RRMHD
convierte ese defecto en un parámetro físico controlado, $\eta=1/\sigma$.}

\Q{nivel básico · lámina 16}{Traduzca sus unidades de código a escalas
astrofísicas.}
\A{Con una longitud de referencia $L_0=10^{16}$~cm (subparsec, típica de la
zona interna de un jet de AGN), el tiempo característico de cruce de luz es
$\tau_c\approx3.8$~días; nuestro horizonte $t=5$ equivale a unos 19 días de
evolución física y el vórtice maduro alcanza $\sim$170 unidades
astronómicas. La física adimensional (números de Mach, Lundquist, $\beta$)
es la que se transporta a otros sistemas.}

\Q{nivel medio · lámina 29}{Su estudio es bidimensional. ¿No invalida eso las
conclusiones?}
\A{No para lo que se mide. La fase lineal del modo fundamental y el
enrollamiento primario son esencialmente bidimensionales; los modos
genuinamente tridimensionales actúan después, fragmentando los vórtices
(Bodo 2004). Lo declaramos como limitación: en 3D esperamos cascada
turbulenta plena y posiblemente otro exponente no lineal. Precisamente por
eso monitoreamos $f_{Vz}$, la fracción de enstrofía fuera del plano, que
señala dónde la dinámica 3D sería más relevante: en la zona de transición.}

\Q{nivel alto · lámina 3}{¿Qué relación tiene la KHI con la aceleración de
partículas que se observa en los jets?}
\A{Indirecta pero central: la KHI genera y adelgaza láminas de corriente en
la capa de mezcla; con resistividad finita esas láminas reconectan y el
campo eléctrico resistivo $\mathbf E'\neq0$ puede acelerar partículas. Es
el escenario invocado para llamaradas rápidas en AGN. Nuestro trabajo
cuantifica el primer eslabón magnetohidrodinámico de esa cadena ---cuánta
corriente y cuánta disipación óhmica produce la inestabilidad según
$\sigma$--- sin modelar las partículas mismas, que exigirían un tratamiento
cinético.}

\section{Marco covariante y geometría}

\Q{nivel básico · lámina 6}{Declare sus convenciones. ¿Por qué es importante
fijarlas una sola vez?}
\A{Unidades de Heaviside--Lorentz con $c=\mu_0=\epsilon_0=1$; signatura
$\eta_{\mu\nu}=\mathrm{diag}(-,+,+,+)$; $W$ es el factor de Lorentz y
$\Gamma$ exclusivamente el índice adiabático (formulación de Valencia,
Rezzolla \& Zanotti 2013). Fijarlas una vez evita errores de signo al
comparar con literatura que usa $(+,-,-,-)$ ---típica de teoría de
campos--- y elimina ambigüedad entre factor de Lorentz e índice adiabático,
que comparten símbolo en parte de la bibliografía.}

\Q{nivel medio · lámina 6 · respaldo B1}{¿Por qué $\nabla\cdot\mathbf B=0$
no es una ecuación dinámica?}
\A{Porque el tensor de Faraday es la derivada exterior del cuadripotencial,
$F=dA$, y la derivada exterior es nilpotente: $dF=d^2A=0$. Esa identidad
---la identidad de Bianchi--- contiene las dos ecuaciones homogéneas de
Maxwell, entre ellas la ausencia de monopolos. Es geometría, no dinámica:
se cumple para cualquier $A$, sin usar las ecuaciones de movimiento. La
dinámica está en las inhomogéneas, $\partial_\mu F^{\mu\nu}=J^\nu$, que sí
provienen de la acción con fuente $J_\mu A^\mu$.}

\Q{nivel alto · respaldo B2}{¿De dónde sale el tensor de energía-momento que
usan? Derívelo conceptualmente.}
\A{Por derivada funcional de la acción respecto de la métrica
(procedimiento de Hilbert): $T^{\mu\nu}=-\frac{2}{\sqrt{-g}}
\frac{\delta S}{\delta g_{\mu\nu}}$. Para el sector electromagnético,
$S_{\rm em}=-\frac14\int F^2\sqrt{-g}\,d^4x$ da el tensor de
Faraday--Maxwell; para el fluido, la acción de Taub $S=\int p\sqrt{-g}$
reproduce el fluido perfecto $\rho h u^\mu u^\nu+p\eta^{\mu\nu}$. El
término de interacción $J_\mu A^\mu$ no contribuye porque no depende de la
métrica. La conservación del total, $\partial_\mu T^{\mu\nu}=0$, expresa la
invariancia bajo traslaciones.}

\Q{nivel medio · lámina 9}{¿Por qué conservan el tensor de energía-momento
\emph{total} y no el del fluido con la fuerza de Lorentz como fuente?}
\A{Ambas formulaciones son equivalentes en el continuo, pero no en la
malla. Al conservar el total, el momento incluye el vector de Poynting y la
energía incluye $\frac12(E^2+B^2)$; la fuerza de Lorentz queda absorbida en
la divergencia del flujo y el esquema de volúmenes finitos conserva energía
y momento \emph{exactamente} a nivel discreto. Con la fuerza como término
fuente, el balance dependería de la cuadratura de la fuente y los choques
se capturarían con saltos incorrectos.}

\Q{nivel medio · lámina 6}{¿Qué papel juega el tensor dual $\star F$?}
\A{Compacta las ecuaciones homogéneas: $\partial_\mu(\star F)^{\mu\nu}=0$
equivale a $dF=0$, es decir, a $\nabla\cdot\mathbf B=0$ y a la ley de
Faraday. En RRMHD el dual es además la vía natural para evolucionar
$\mathbf B$: el sector magnético del sistema conservativo es exactamente la
componente espacial de esa ecuación. En la formulación resistiva, a
diferencia de la ideal, $\mathbf E$ es variable independiente y las
inhomogéneas se evolucionan explícitamente con la corriente de Ohm.}

\Q{nivel alto · lámina 6}{Muestre que su formulación recupera la fuerza de
Lorentz clásica en el límite no relativista.}
\A{La divergencia espacial de $T^{\mu\nu}_{\rm em}$ da, con $W\to1$ y
$v\ll c$, $\partial_t(\mathbf E\times\mathbf B)+\nabla\cdot\mathsf T=
-(\rho_q\mathbf E+\mathbf J\times\mathbf B)$, con $\mathsf T$ el tensor de
esfuerzos de Maxwell: el intercambio de momento entre campo y materia es
exactamente $\rho_q\mathbf E+\mathbf J\times\mathbf B$, la fuerza de
Lorentz de la electrodinámica de pregrado. En el balance del fluido
aparece con signo opuesto, como debe ser por tercera ley.}

\Q{nivel básico · lámina 9}{Enumere sus 14 variables conservadas y por qué
son 14.}
\A{Densidad relativista $D=\rho W$; tres momentos totales $S^j$ (incluyen
Poynting); energía total $\tau$; tres componentes de $\mathbf B$; tres de
$\mathbf E$ ---independientes porque el cierre es resistivo---; la densidad
de carga $\rho_q$; y los dos escalares de limpieza $\Psi,\Phi$ del sistema
GLM. Total: $1+3+1+3+3+1+2=14$. En RMHD ideal serían 8, porque $\mathbf E$
deja de ser independiente y no se evoluciona la carga ni la limpieza
eléctrica.}

\Q{nivel medio · respaldo B5}{¿Qué diferencia hay entre $\rho_e$ y $\rho_q$
en su ley de Ohm?}
\A{$\rho_e$ es la densidad de carga medida en el marco comóvil del fluido
---la que multiplica a $u^\mu$ en la forma covariante---; $\rho_q$ es la
densidad de carga en el marco del laboratorio, la que se evoluciona en la
malla euleriana. Se relacionan por la proyección $3+1$:
$\rho_q=\rho_e W+\sigma W\,(\mathbf v\cdot\mathbf E)/c$ en nuestra
normalización. Confundirlas produce una ley de Ohm inconsistente con la
conservación de la carga.}

\Q{nivel alto · lámina 6}{¿Su sistema es invariante gauge? ¿Dónde vive el
cuadripotencial en el código?}
\A{La física lo es: el código no evoluciona $A^\mu$ sino los campos
$F^{\mu\nu}$ ---$\mathbf E$ y $\mathbf B$---, que son invariantes gauge. El
potencial solo interviene conceptualmente, para demostrar que $dF=0$ es
identidad (respaldo B1) y en el origen variacional del GLM, donde los
escalares de limpieza actúan como campos auxiliares tipo Stueckelberg que
restauran la consistencia de las ligaduras (Lee 2004, respaldo B3).}

\section{Ley de Ohm, resistividad y fundamento cinético}

\Q{nivel básico · lámina 7}{Escriba la ley de Ohm relativista y señale su
huella relativista.}
\A{Covariante: $J^\mu=\rho_e u^\mu+\sigma F^{\mu\nu}u_\nu$, respuesta
lineal e isótropa en el marco comóvil. Proyectada al laboratorio:
$\mathbf J=\rho_q\mathbf v+\sigma W[\mathbf E+\mathbf v\times\mathbf B
-(\mathbf E\cdot\mathbf v)\mathbf v]$. La huella relativista es el factor
$W$ y el término $(\mathbf E\cdot\mathbf v)\mathbf v$: no es la suma
galileana de conducción más convección, sino la transformación de Lorentz
del campo eléctrico comóvil.}

\Q{nivel básico · lámina 7}{¿Qué ocurre en los límites $\sigma\to\infty$ y
$\sigma\to0$?}
\A{Con $\sigma\to\infty$, para que la corriente sea finita debe anularse el
campo comóvil: $\mathbf E+\mathbf v\times\mathbf B=0$, el congelamiento
ideal. Con $\sigma\to0$ la materia y el campo se desacoplan: quedan las
ecuaciones de Maxwell en vacío (con la carga advectada) y la hidrodinámica
pura. Entre ambos, el campo comóvil sobrevive y produce calentamiento de
Joule ($j^\mu E_\mu>0$) y reconexión: el régimen de este trabajo.}

\Q{nivel alto · lámina 7 · respaldo B9}{¿Cuál es el origen microscópico de
la resistividad que usted puso a mano?}
\A{Es un coeficiente de transporte que emerge al truncar la jerarquía
cinética: partiendo del teorema de Liouville, la ecuación de
Vlasov--Boltzmann para la función de distribución incluye integrales de
colisión; al tomar momentos y cerrar en el nivel unifluido, la
transferencia de momento entre especies por colisiones se condensa en
$\eta$. El cierre escalar e isótropo exige colisionalidad alta
($\nu_{\rm col}\gg\omega_{\rm cicl}$), que borra la memoria direccional del
campo. En plasmas reales $\eta$ depende de la temperatura (tipo Spitzer);
aquí se toma constante y uniforme para aislar su efecto dinámico como
parámetro de control.}

\Q{nivel medio · lámina 7}{¿Por qué es lícito despreciar el efecto Hall y la
anisotropía de $\sigma$?}
\A{Porque el régimen declarado es colisional y macroscópico: si la
frecuencia de colisiones domina sobre la de ciclotrón, los electrones no
completan giros entre colisiones y la conducción no distingue direcciones
respecto de $\mathbf B$; el término de Hall, del orden del cociente entre
la escala iónica y la escala del sistema, es despreciable cuando la escala
disipativa relevante es la capa resistiva $\delta\sim S^{-1/2}L$ y no el
girorradio. Lo enunciamos como límite de validez, no como teorema: fuera de
él se necesitaría Ohm generalizada o descripción cinética.}

\Q{nivel medio · lámina 23}{Defina el número de Lundquist y ubique su
transición en esa escala.}
\A{$S=\tau_\eta/\tau_A=Lv_A/\eta$: cociente entre el tiempo de difusión
resistiva y el de cruce de Alfvén. Con las escalas de nuestra capa, la zona
de transición $\sigma\in[1400,7000]$ corresponde a $S\sim10^3$. Es un valor
físicamente razonable: por debajo la difusión domina la dinámica de la
capa; por encima el congelamiento es efectivo durante la vida de la
inestabilidad, y coincide con el orden en que la literatura de reconexión
sitúa la activación de dinámicas de láminas inestables.}

\Q{nivel básico · lámina 14}{¿Por qué su sistema es rígido y qué escala lo
gobierna?}
\A{El término fuente de Ampère es $-\sigma\mathbf E$ (más los acoples de
velocidad): induce una relajación del campo eléctrico hacia su valor de
Ohm en un tiempo $\tau_{\rm relax}\sim1/\sigma$. A $\sigma=10^4$ ese tiempo
es cuatro órdenes de magnitud menor que el paso advectivo permitido por
Courant. Un integrador explícito debería resolverlo ($\Delta t\lesssim
1/\sigma$) aunque la física de interés viva en el tiempo dinámico: esa
disparidad es la rigidez.}

\Q{nivel medio · lámina 26}{¿Dónde aparece $\eta$ en el balance energético y
con qué signo?}
\A{En el trabajo del campo sobre la materia, $\mathbf E\cdot\mathbf J$.
Separando la corriente de Ohm, el término conductivo da
$\sigma W|\mathbf E'|^2\ge0$ en el comóvil: disipación de Joule
estrictamente positiva, irreversible, que convierte energía
electromagnética en interna. Por eso la integral $\int\mathbf E\cdot
\mathbf J\,dA$ es nuestro diagnóstico de disipación óhmica global, y por
eso la resistividad decide si el intercambio cinético-magnético de la
campaña B es reversible ($\sigma=10^4$) o termina en calor ($\sigma=6000$).}

\section{Divergencias, GLM y estructura del sistema aumentado}

\Q{nivel básico · lámina 8}{Si $\nabla\cdot\mathbf B=0$ es una identidad,
¿por qué su código tiene que ``limpiarla''?}
\A{Porque la identidad depende de la nilpotencia $d^2=0$, y los operadores
en diferencias de una malla no son nilpotentes: el rotacional discreto de
un gradiente discreto no es exactamente cero. El error de truncamiento
inyecta una divergencia espuria ---monopolos numéricos--- que, sin
tratamiento, se acumula, produce fuerzas paralelas artificiales a
$\mathbf B$ y desestabiliza la simulación (Dedner 2002).}

\Q{nivel medio · lámina 8 · respaldo B3}{Explique el mecanismo GLM y por qué
no destruye la causalidad.}
\A{Se agregan dos escalares $\Phi,\Psi$ ---multiplicadores de Lagrange
generalizados--- que acoplan las ligaduras de Gauss magnética y eléctrica a
la evolución. Eliminando los multiplicadores, el error de divergencia
obedece la ecuación del telégrafo:
$\partial_t^2\xi+\kappa\partial_t\xi-c_h^2\nabla^2\xi=0$. Es hiperbólica
---el error se \emph{propaga} con velocidad $c_h=1$, dentro del cono de
luz--- y disipativa ---decae como $e^{-\kappa t/2}$---. El error ni se
acumula ni viaja superlumínicamente: se evacúa y se amortigua.}

\Q{nivel alto · respaldo B3}{¿Cuál es el origen variacional del GLM?}
\A{Se obtiene extendiendo la acción electromagnética con términos tipo
Stueckelberg en los escalares de limpieza, que restauran como dinámica lo
que la discretización rompió como identidad (Munz 2000 para Maxwell; Lee
2004 para la estructura lagrangiana y de simetrías). La forma de segundo
orden (telégrafo) es equivalente a la formulación de primer orden de
Dedner que integra el código; la ventaja de la primera es conceptual
---exhibe causalidad y disipación--- y la de la segunda es práctica: encaja
en el esquema hiperbólico de volúmenes finitos.}

\Q{nivel alto · lámina 9}{¿Por qué su vector fuente contiene términos de
Godunov--Powell si son no conservativos?}
\A{Porque mientras el error de divergencia es transitoriamente distinto de
cero, la forma puramente conservativa acelera los monopolos numéricos y
viola la compatibilidad termodinámica. Los términos proporcionales a
$\Psi$ ---$-\Psi B^j$ en el momento y $-\Psi(\mathbf v\cdot\mathbf B)$ en
la energía--- compensan exactamente ese defecto y hacen el sistema
compatible con una desigualdad de entropía (Rueda-Ramírez 2023). Son
formalmente nulos en el continuo, donde $\nabla\cdot\mathbf B=0$: no
cambian la física, blindan la termodinámica del transitorio numérico.}

\Q{nivel medio · lámina 8}{¿Monitorearon el error de divergencia? ¿Qué
comportamiento tiene?}
\A{Sí: la norma del error se registra en el tiempo. Con GLM activo
permanece acotada y decae tras cada transitorio; la figura de la lámina 8
contrasta la acumulación sin tratamiento con la propagación y el
amortiguamiento con GLM. Además la ligadura eléctrica ---ley de Gauss con
$\rho_q$--- se trata con el segundo escalar, de modo que ambas divergencias
quedan controladas con el mismo mecanismo telegráfico.}

\Q{nivel medio · lámina 8}{¿Por qué GLM y no transporte restringido
(constrained transport)?}
\A{CT mantiene $\nabla\cdot\mathbf B$ a error de máquina, pero exige mallas
escalonadas y una discretización específica del rotacional que complica el
acople con reconstrucción de quinto orden, con el sector eléctrico
independiente de la RRMHD y con la recuperación de primitivas. GLM
mantiene todas las variables colocalizadas, encaja en el mismo esquema
hiperbólico de las demás ecuaciones, controla también la ligadura
eléctrica y su error es acotado, monitoreado y físicamente inocuo. Es la
solución de compromiso estándar en los códigos RRMHD de referencia
(Komissarov 2007; Palenzuela 2009; Miranda-Aranguren 2018).}

\section{Hiperbolicidad, espectro característico y límites}

\Q{nivel medio · respaldo B4}{Describa el espectro del jacobiano de su
sistema.}
\A{El jacobiano $14\times14$ de los flujos tiene 14 autovalores reales: 8
degenerados en $\pm c$ ---los sectores electromagnético y GLM propagan a
la velocidad de la luz---, las magnetosónicas rápidas del sector fluido,
los modos de entropía y cizalla que viajan con el fluido, y el modo de
advección de carga. El sistema es hiperbólico; la degeneración en $\pm c$
es la firma de que Maxwell no se ``esclaviza'' al fluido como en RMHD
ideal.}

\Q{nivel alto · respaldo B4}{¿Qué garantiza la estabilidad del sistema
relajado frente al rígido? (condición subcaracterística)}
\A{La condición de Liu (1987): en un sistema hiperbólico de relajación, las
velocidades características del sistema de equilibrio ---aquí la RMHD
ideal que emerge cuando $\sigma\to\infty$--- deben quedar entrelazadas por
las del sistema completo. Se cumple: las magnetosónicas ideales son
subluminales y el sistema resistivo propaga hasta $\pm c$. Esto asegura
que el límite rígido no produce inestabilidades de relajación y da soporte
teórico a la preservación asintótica del integrador.}

\Q{nivel básico · lámina 17}{¿Cómo fijan el paso de tiempo? ¿Por qué CFL
distintos entre campañas?}
\A{Por la condición de Courant sobre la velocidad característica máxima,
que en RRMHD es $c$: $\Delta t=\mathrm{CFL}\cdot\Delta x/c$. El barrido usa
CFL $=0.1$; las campañas magnéticas a $\sigma=6000$ y $10000$ usan $0.04$ y
$0.02$ porque el campo fuerte y las láminas de corriente agudas estresan la
recuperación de primitivas, y reducir el paso da margen de robustez. La
rigidez resistiva no limita $\Delta t$: la absorbe la parte implícita del
IMEX.}

\Q{nivel medio · lámina 10}{¿Por qué no existe relación de dispersión
cerrada en RRMHD y qué consecuencia metodológica tiene?}
\A{El análisis lineal clásico de la discontinuidad tangencial empalma
soluciones a ambos lados de una interfaz que se supone persistente. En
RRMHD el término $\eta\nabla^2\mathbf B$ es parabólico: difunde la
discontinuidad en tiempo arbitrariamente corto, de modo que el estado base
del análisis modal deja de existir y el empalme algebraico es inválido
(Palenzuela 2009). Consecuencia: $\gamma_{\rm KHI}(\sigma)$ debe
caracterizarse numéricamente, con la teoría lineal ideal como referencia
asintótica. Esa imposibilidad es la justificación epistemológica del
trabajo.}

\Q{nivel básico · lámina 10}{Recorra el mapa de límites de su sistema.}
\A{RRMHD con $\sigma\to\infty$ da RMHD ideal ($\mathbf E'=0$, congelamiento);
con $v\ll c$ y $W\to1$, MHD clásica; con $\mathbf B\to0$, hidrodinámica de
Euler; y con $\sigma\to0$, Maxwell en vacío más hidrodinámica desacoplada.
En paralelo, la KHI: incondicionalmente inestable en hidrodinámica; en MHD
la tensión estabiliza si $(\mathbf k\cdot\mathbf B)^2$ basta; en RMHD rige
$0<M_{re}<\sqrt2$; y en RRMHD la tasa se vuelve función medible de
$\sigma$, que es exactamente lo que medimos.}

\Q{nivel alto · lámina 14}{¿En qué sentido su esquema ``respeta el límite
ideal por construcción''?}
\A{Preservación asintótica: en la actualización implícita del campo
eléctrico, al dividir por $\alpha\sigma W$ y tomar $\sigma\to\infty$, la
solución converge algebraicamente a $\mathbf E=-\mathbf v\times\mathbf B$,
la ley de Ohm ideal, sin pasos intermedios inestables ni necesidad de
resolver la escala $1/\sigma$. El esquema discreto reproduce el mismo
límite que el sistema continuo; por eso pudimos barrer hasta
$\sigma=10^4$ con el paso de tiempo dictado solo por la advección.}

%========================= PARTE B: categorías 6–10 ==========================

\section{Métodos numéricos}

\Q{nivel básico · lámina 12}{¿Por qué volúmenes finitos y no elementos
finitos?}
\A{Por el tipo de solución. FEM, con su formulación débil global, es óptimo
para dominios complejos con soluciones suaves. En flujos relativistas
emergen choques que vuelven singular la forma diferencial: se necesitan
soluciones débiles con saltos, y la formulación integral de FVM las admite
por construcción; los flujos numéricos de Riemann respetan las condiciones
de salto de Rankine--Hugoniot (Toro 2009). Además la conservación discreta
es exacta: el flujo que sale de una celda entra en la vecina, requisito
para capturar choques con las velocidades correctas.}

\Q{nivel medio · lámina 13}{Enuncie el teorema de Godunov y explique cómo lo
esquiva MP5.}
\A{Ningún esquema lineal de orden mayor que uno preserva monotonicidad:
alto orden lineal implica oscilaciones de Gibbs en discontinuidades. La
salida es la no linealidad del esquema: limitadores. Los TVD de segundo
orden (minmod y afines) son monótonos pero recortan los extremos suaves
---y con ello destruyen la vorticidad que queremos medir---. MP5 (Suresh \&
Huynh 1997) usa un \emph{stencil} de cinco puntos con límite de
preservación de monotonicidad mediante una mediana proyectiva: mantiene
quinto orden en regiones suaves, no oscila en choques y retiene los
extremos de vorticidad de forma casi espectral.}

\Q{nivel básico · lámina 13}{¿Qué es el flujo HLL y cuáles son sus
propiedades?}
\A{Un solucionador de Riemann aproximado de dos ondas: supone un único
estado intermedio entre las velocidades características extremas
$\lambda_L,\lambda_R$ y da el flujo
$\mathbf F^{hll}=[\lambda_R\mathbf F_L-\lambda_L\mathbf F_R+
\lambda_L\lambda_R(\mathbf U_R-\mathbf U_L)]/(\lambda_R-\lambda_L)$
(Harten, Lax \& van Leer 1983). Es robusto, positivo, libre de jacobianos
---valioso con un sistema de 14 variables--- y entrópico; su precio es
difusión sobre la onda de contacto, que en nuestro esquema compensa la
agudeza de los estados reconstruidos a quinto orden.}

\Q{nivel alto · respaldo B7}{El artículo del código presenta un solucionador
HLLC. ¿Por qué el barrido usa HLL?}
\A{Decisión de diseño documentada. HLLC restaura la onda de contacto
explícitamente; pero Miranda-Aranguren et al.\ (2018) muestran que HLL
alimentado con estados MP5 alcanza una nitidez de contacto comparable: la
información que HLLC añade en el solucionador, MP5 la aporta en la
reconstrucción. Para un barrido de 62 simulaciones elegimos la combinación
más robusta (positividad garantizada, sin estados intermedios delicados) y
más barata por paso. La verificación de que la difusión numérica resultante
quedó por debajo de la física es empírica: el plateau de $J_{\max}$.}

\Q{nivel alto · láminas 13 y 26}{¿Cómo garantiza que su vórtice es física y
no viscosidad artificial?}
\A{En tres niveles. Diseño: MP5$+$HLL se eligió para que la disipación
numérica quede por debajo de la física en todo el barrido. Evidencia
empírica: el plateau de $J_{\max}$ localiza dónde la difusividad física
$\eta$ supera el error de truncamiento; la zona de transición completa está
en ese régimen. Conceptual: la KHI 2D \emph{ideal} ni siquiera converge con
la resolución (Lecoanet 2016), de modo que es precisamente la $\eta$
explícita la que da significado convergente y bien planteado a
$\gamma_{\rm KHI}(\sigma)$.}

\Q{nivel medio · lámina 14}{Describa la partición IMEX y el costo real de su
parte implícita.}
\A{Los flujos hiperbólicos se integran explícitamente (MP5$+$HLL, $\Delta t
\propto\Delta x$); la fuente rígida $-\sigma\mathbf E$, implícitamente, en
un Runge--Kutta IMEX de tipo SSP (Pareschi \& Russo 2005). Como la fuente
no acopla celdas vecinas, la etapa implícita es un sistema algebraico local
$3\times3$ para $\mathbf E$ con solución analítica: no hay Newton global ni
inversión iterativa, y el sobrecosto por paso es despreciable frente al
salto de paso temporal que se gana ---de $\Delta t\propto\Delta x^2/\sigma$
a $\Delta t\propto\Delta x$.}

\Q{nivel medio · lámina 15}{¿Cómo recupera el código las variables
primitivas y qué guardias tiene?}
\A{El mapeo conservadas$\to$primitivas no tiene inversa cerrada porque $W$
acopla cinemática y termodinámica. Cueva lo reduce, con la EoS politrópica,
a una cuártica en $W$ que se resuelve analíticamente (Tschirnhaus $+$
Ferrari--Cardano) y se refina con pocas iteraciones de Newton--Raphson.
Guardias: rechazo de $v^2\ge1$ y de $p\le0$, rama linealizada estable a
bajas velocidades y \emph{fallback} a promedios de celdas vecinas si la
raíz es inadmisible. Se ejecuta en cada celda y subpaso: su robustez define
el límite práctico del barrido.}

\Q{nivel medio · lámina 17}{¿Qué recursos computacionales consumió el
trabajo y en qué arquitectura?}
\A{Unas $3\times10^4$ horas-núcleo acumuladas en CPU. Cueva es Fortran
paralelo sobre CPU ---no CUDA---, con volúmenes finitos unigrid; el
análisis posterior es Python. La elección es deliberada: para mallas
$512\times256$ en 2D el cuello es el barrido de parámetros, no la
simulación individual, y la robustez portable de CPU simplifica reproducir
cualquier punto del barrido.}

\Q{nivel alto · lámina 19}{¿Hicieron estudio de convergencia de malla? ¿Qué
responde a la exigencia de independencia de malla?}
\A{Con matiz: para la KHI 2D ideal el estudio de convergencia clásico está
condenado ---Lecoanet et al.\ (2016) demuestran que sin disipación
explícita las cantidades no convergen al refinar---. Nuestro planteamiento
invierte la lógica: la $\eta$ física fija la escala disipativa, y la
evidencia de resolución suficiente es que la escala resistiva esté resuelta
donde extraemos conclusiones cuantitativas: el plateau de $J_{\max}$ y las
$\approx13$ celdas por capa de cizalla lo acreditan. Declaramos que a las
$\sigma$ más altas las láminas rozan la malla, y por eso el exponente
$1.22$ se presenta con esa salvedad.}

\section{Montaje experimental y diseño de campañas}

\Q{nivel básico · lámina 16}{Describa su condición inicial. ¿Qué gatilla la
inestabilidad?}
\A{Doble capa de cizalla tipo Mizuno (2013): perfiles
$V_y(x)=\pm v_{\rm sh}\tanh[(x\mp0.5)/a_{kh}]$ con $v_{\rm sh}=0.5$ y
$a_{kh}=0.05$; contraste de densidad $\eta_\rho=0.1$ (chorro diluido en
ambiente denso); campo $\mathbf B=(0,\sqrt{0.02},\sqrt{2.0})$;
$\mathbf E(0)=0$. La semilla es una perturbación senoidal en $V_x$, el modo
fundamental $k=2\pi/L_y$ con confinamiento gaussiano en las capas. El
contraste de densidad modela la configuración chorro--medio pero no gatilla
nada: sin la perturbación de velocidad el estado base es estacionario.}

\Q{nivel medio · lámina 16}{¿Por qué eligieron un campo guía $B_z$
dominante?}
\A{Es la decisión topológica que hace posible el experimento: con
$\mathbf k$ en $\hat y$ y el campo casi todo en $\hat z$, el producto
$\mathbf k\cdot\mathbf B\propto B_y$ es débil, de modo que el campo aporta
presión magnética pero casi ninguna tensión. Un campo fuerte paralelo a
$\mathbf k$ habría estabilizado la KHI y no habría nada que medir. Así la
inestabilidad se desarrolla y el efecto de $\sigma$ queda aislado; la
componente $B_y=\sqrt{0.02}$ pequeña conserva la semilla de dínamo que
señala Mizuno.}

\Q{nivel medio · lámina 16 · respaldo B6}{¿Las dos capas de cizalla no se
contaminan mutuamente? ¿Y las fronteras?}
\A{No: el acople entre capas decae como $e^{-kD}$ con la separación $D$;
para nuestra geometría es del orden del $0.2\%$, verificado con el
solucionador lineal (respaldo B6). Las condiciones de frontera son
periódicas en $y$ ---compatibles con el modo fundamental--- y abiertas en
$x$, con el dominio suficientemente ancho para que las ondas emitidas
salgan sin reflejarse sobre la capa durante el horizonte simulado.}

\Q{nivel básico · lámina 16}{¿Por qué $\Gamma=4/3$?}
\A{Es el índice adiabático del gas relativista caliente (fotones o pares, o
bariones con presión dominada por especies relativistas), el régimen
apropiado para el interior de chorros de AGN. La EoS es politrópica por
diseño: Mizuno (2013) muestra que la elección de EoS modula la KHI
resistiva, así que fijarla elimina un eje de variación y deja $\sigma$ y
$\mathbf B$ como únicos parámetros de control.}

\Q{nivel básico · lámina 17}{Presente su matriz de simulaciones.}
\A{Campaña 1: barrido de 44 conductividades $\sigma\in[100,10\,500]$ con
CFL $0.1$, todo lo demás fijo. Campaña 2: a $\sigma=6000$ (CFL $0.04$) y
$\sigma=10\,000$ (CFL $0.02$), tres familias: A ---intensidad $B_0=1.0,
1.5,2.0$---, B ---rotación del campo en el plano $y$--$z$, que activa
tensión--- y C ---rotación en $x$--$z$ con $\mathbf k\cdot\mathbf B=0$, sin
tensión---. Todas con malla $512\times256$ y horizonte $t>5$.}

\Q{nivel medio · lámina 17}{¿Qué datos excluyeron y con qué criterio?}
\A{Tres exclusiones, declaradas: en la campaña A, $B_0=0.25$ y $0.5$
fallaron por inestabilidad numérica antes de la fase no lineal; en la B,
$\theta=15^\circ$ se detuvo por fallo de convergencia del solucionador de
primitivas; y en el barrido, el análisis cuantitativo de $\gamma$ usa
$\sigma\ge1600$ con $R^2\ge0.90$ y señal-ruido $\ge3$ ---por debajo no
existe fase lineal que ajustar, no es un ajuste pobre---. Los casos
excluidos se conservan y se reportan: la transparencia es parte del
resultado.}

\Q{nivel medio · lámina 16}{¿Es consistente iniciar con $\mathbf E=0$ si la
ley de Ohm exige otro campo?}
\A{Sí, y es deliberado: el término rígido $-\sigma(\mathbf E-
\mathbf E_{\rm ohm})$ relaja el campo eléctrico hacia el valor consistente
con la conductividad en un tiempo $\sim1/\sigma$, órdenes de magnitud más
corto que el crecimiento de la KHI. El transitorio inicial queda fuera de
la ventana de ajuste $t\in[2.4,3.4]$ precisamente para que esta relajación
---y el acomodo acústico inicial--- no contaminen la medición.}

\Q{nivel alto · lámina 17}{¿Por qué el rango $[100,10\,500]$ y no
conductividades mayores?}
\A{El extremo inferior garantiza cubrir el régimen dominado por difusión; el
superior está dictado por la resolución: a $\sigma$ mayores la capa
resistiva $\delta\sim S^{-1/2}$ cae por debajo de la malla y los
resultados dejarían de estar en el régimen controlado que exigimos
(disipación física $>$ numérica). Con $512\times256$, $\sigma=10\,500$ es
la frontera prudente; medir el plateau ideal directamente exige
$\sigma\sim2$--$3\times10^4$ con más resolución, y así se declara como
trabajo futuro.}

\section{Observables y metodología de análisis}

\Q{nivel básico · lámina 11}{Defina su observable primario.}
\A{La enstrofía de perturbación. A la vorticidad $\omega_z=\partial_xV_y-
\partial_yV_x$ se le sustrae el promedio en $y$ de la configuración
inicial, para remover la cizalla base: $\omega_z'=\omega_z-\langle\omega_z
(t{=}0)\rangle_y$. Se integra su cuadrado: $\Omega_{zp}=\int(\omega_z')^2
dA$. Mide la intensidad rotacional de la perturbación, es positiva,
integral (robusta al ruido puntual) y estándar en la literatura de KHI
resistiva (Nastro 2022).}

\Q{nivel básico · lámina 11}{¿Por qué la pendiente de $\ln\Omega_{zp}$ es
$2\gamma$ y no $\gamma$?}
\A{Porque en fase lineal cada modo crece como $e^{\gamma t}$ y la enstrofía
es \emph{cuadrática} en la amplitud de la perturbación: su logaritmo crece
con el doble de la tasa. Hacerlo explícito evita el error de factor 2 más
común al comparar tasas entre trabajos que usan observables lineales
(amplitud de modo) y cuadráticos (energías, enstrofías).}

\Q{nivel medio · lámina 20}{¿Cómo eligieron la ventana de ajuste
$t\in[2.4,3.4]$?}
\A{Es posterior al transitorio inicial ---relajación de Ohm y acomodo
acústico de la condición inicial--- y anterior a la saturación del modo
fundamental incluso en los casos más rápidos. Se usa la \emph{misma}
ventana para las 44 simulaciones: una ventana fija sacrifica algo de ajuste
en los casos extremos, pero garantiza que las tasas sean comparables entre
sí y no un artefacto de ventanas elegidas a conveniencia.}

\Q{nivel básico · lámina 11}{Enumere sus diagnósticos complementarios y qué
mide cada uno.}
\A{$\Omega_{\rm tot}$ y la fracción fuera del plano $f_{Vz}=(\Omega_{\rm
tot}-\Omega_z)/\Omega_{\rm tot}$: estructuras secundarias con velocidad
$V_z$. $J_{\max}$: el adelgazamiento de las láminas de corriente.
$\int\mathbf E\cdot\mathbf J\,dA$: disipación óhmica global. Las energías
$E_{\rm kin}$, $E_{\rm mag}$, $E_{\rm int}$ y su suma: canales de
conversión y control de conservación. Y la tasa normalizada
$\hat\omega=\gamma\,a_{kh}/v_{\rm sh}$ para comparar con la literatura.}

\Q{nivel medio · lámina 11}{¿Por qué enstrofía y no la amplitud de un modo de
Fourier?}
\A{La amplitud modal es óptima para un problema estrictamente lineal, pero
aquí el fondo se difunde (el perfil $\tanh$ se ensancha con $\eta$), lo que
contamina los modos bajos, y en el régimen resistivo la energía se
redistribuye entre armónicos desde tiempos tempranos. La enstrofía de
perturbación integra toda la actividad rotacional de la perturbación, es
insensible a esa redistribución y sigue siendo válida cuando la fase lineal
es marginal ---exactamente el régimen que queríamos clasificar.}

\Q{nivel medio · lámina 26}{¿Qué significa físicamente el máximo de
$f_{Vz}$ en la transición?}
\A{$f_{Vz}$ mide cuánta enstrofía vive en movimientos fuera del plano. En el
régimen resistivo la difusión suprime los gradientes que alimentarían
$V_z$; en el cuasi-ideal el congelamiento rigidiza el plasma y también las
suprime. El máximo ($f_{Vz}=0.37$ en $\sigma\approx1400$) aparece donde
difusión y advección compiten: la zona de transición es el hábitat de las
estructuras secundarias. Es uno de los cuatro relojes que usamos para
delimitar la zona crítica.}

\section{Resultados del barrido: fase lineal y transición}

\Q{nivel básico · lámina 21}{Presente su resultado paramétrico central.}
\A{La tasa lineal en función de la conductividad queda descrita por un
sigmoide con desplazamiento:
$\gamma_{\rm KHI}(\sigma)=C+\gamma_0/(1+e^{-k(\sigma-\sigma_0)})$. La
asíntota ideal vale $C+\gamma_0=1.03\pm0.05$ (normalizada,
$\hat\omega=0.103$) y el centro de la transición es
$\sigma_0=2815\pm123$. El modelo se ajusta sobre las 35 simulaciones con
fase lineal medible y supera decisivamente al sigmoide sin desplazamiento.}

\Q{nivel medio · lámina 21}{El desplazamiento $C$ del ajuste es negativo.
¿Una tasa de crecimiento negativa?}
\A{No: $C$ aislado no es un observable físico, es un parámetro del modelo
---la extrapolación formal del sigmoide hacia $\sigma\to-\infty$, fuera del
dominio físico---. Lo defendible es su papel funcional: con $C$ libre el
modelo predice un 36\% mejor los nueve puntos resistivos excluidos del
ajuste. Además $C$ y $\gamma_0$ están fuertemente anticorrelacionados
($r=-0.997$); la combinación robusta es la suma $C+\gamma_0$, que varía
menos del 4\% entre variantes del ajuste y es la que interpretamos.}

\Q{nivel medio · lámina 22}{¿Cómo estimaron $\sigma_{\rm crit}$ y por qué
cuatro veces?}
\A{Con un método común aplicado a cuatro observables independientes: el
punto de máxima derivada respecto de $\log_{10}\sigma$ (inflexión), con
error el máximo entre el remuestreo jackknife y el semiancho de la zona de
inflexión. Resultados: $2815\pm123$ (sigmoide de $\gamma$), $3400\pm2175$
($\Omega^{\max}_{zp}$), $6000\pm1147$ ($t_{\rm peak}$) y $6800\pm2625$
($\gamma$ en escala log). Cuatro relojes distintos porque cada fenómeno se
``idealiza'' a su propia conductividad: esa es la física del resultado.}

\Q{nivel básico · lámina 23}{¿Por qué insisten en que $\sigma_{\rm crit}$ es
una zona y no un número?}
\A{Porque la dispersión entre estimadores ($\pm1944$) excede en un orden de
magnitud el error formal del mejor ajuste ($\pm123$): reportar
``$2815\pm123$'' sería precisión espuria. La lectura honesta es una zona de
transición $\sigma\in[1400,7000]$ ---del máximo de $f_{Vz}$ a la
saturación no lineal--- con centro $\approx2815$. La dispersión no es
ruido: es la firma de una transición multiescala en la que cada observable
se congela a su escala. En número de Lundquist, $S\sim10^3$.}

\Q{nivel medio · lámina 20}{¿Cómo clasificaron los regímenes del barrido?}
\A{Por la calidad del ajuste lineal en la ventana canónica: $R^2<0.90$
define el régimen resistivo (no hay recta que ajustar), $0.90$--$0.995$ la
transición y $\ge0.995$ el cuasi-ideal. Es un criterio operativo y
reproducible que coincide con la morfología: difusión de vórtices en el
resistivo, subestructura fina y turbulencia de pequeña escala en el
cuasi-ideal.}

\Q{nivel medio · lámina 20}{Para $\sigma\lesssim1400$ no reportan tasa.
¿Fracasó el método?}
\A{No: el resultado es que \emph{no existe} fase lineal que medir. Con
difusividad grande, la capa se ensancha en el tiempo de crecimiento del
modo: la difusión domina desde $t=0$ y el crecimiento exponencial nunca se
establece. Forzar una pendiente daría un número sin significado. Esas
simulaciones se conservan para la clasificación de regímenes y para los
diagnósticos globales, que no requieren fase lineal.}

\Q{nivel básico · lámina 19}{Describa las tres fases temporales que exhibe
el barrido.}
\A{Fase lineal hacia $t\sim3$: la perturbación crece exponencialmente y la
interfaz se ondula. Enrollamiento no lineal hacia $t\sim8$: los vórtices
maduran, generan láminas de corriente y, según $\sigma$, reconectan o
enrollan el campo. Saturación hacia $t\sim14$: la enstrofía alcanza su
máximo y decae mientras la capa se ensancha. Las 44 simulaciones muestran
la secuencia; lo que cambia con $\sigma$ es la tasa, la morfología y el
destino de la energía.}

\Q{nivel alto · lámina 21}{¿Qué garantiza que el sigmoide no es un
sobreajuste con 4 parámetros?}
\A{Tres controles. Comparación anidada: contra el sigmoide de 3 parámetros
($C=0$), sobre los mismos datos, $\Delta\mathrm{AIC}=-126$; en la escala de
Burnham \& Anderson, evidencia decisiva que paga con creces el parámetro
extra. Validación fuera de muestra: mejora del 36\% (RMSE
$0.039\to0.025$) sobre los 9 puntos resistivos no usados. Y honestidad
residual: los residuos están autocorrelacionados, así que no confiamos el
error a la matriz de covarianza del ajuste sino a la dispersión
inter-método de la lámina siguiente.}

\section{Contraste con la teoría lineal}

\Q{nivel básico · lámina 24}{Explique la ``escalera de predicciones''.}
\A{Tres peldaños. Techo hidrodinámico inviscido: la ecuación de Rayleigh
para el perfil $\tanh$ da $\hat\omega_{\max}=0.190$ (Michalke 1964).
Predicción magnetosónica: como $\mathbf k\perp\mathbf B$, el campo aporta
presión pero no tensión y la velocidad de restitución es
$v_{f\perp}=0.691$; la teoría compresible tipo Lees--Lin da entonces
$\hat\omega=0.095$. Valor medido: la asíntota del sigmoide,
$\hat\omega=0.103$. Acuerdo del 8\% sin parámetros de ajuste.}

\Q{nivel medio · lámina 24 · respaldo B6}{¿Cómo validaron su solucionador de
teoría lineal?}
\A{Reproduciendo el resultado clásico: para el perfil $\tanh$ inviscido e
incompresible, el solucionador de Rayleigh recupera
$\hat\omega_{\max}=0.1897$ de Michalke (1964) en el número de onda
correcto. Además verifica el umbral compresible ($\sqrt2\,c_s$ del caso
\emph{vortex sheet}, Landau 1944) y el desacople de las dos capas
($e^{-kD}\approx0.2\%$). Sobre esa base validada se aplica la
sustitución magnetosónica.}

\Q{nivel medio · lámina 24}{Justifique la sustitución $c_s\to v_{f\perp}$.}
\A{Es el resultado de Chow, Rosen \& Piran (2023) para geometría de campo
guía: cuando $\mathbf k\perp\mathbf B$, las perturbaciones comprimen el
campo sin doblarlo; el campo no ejerce tensión restauradora sobre el modo,
pero su presión sí rigidiza el medio. La velocidad efectiva de restitución
pasa de $c_s$ a la magnetosónica rápida perpendicular
$v_{f\perp}=\sqrt{c_s^2+v_A^2}$ relativista, $0.691$ en nuestro estado
base. ``Presión pero no tensión'' resume la física de la sustitución.}

\Q{nivel medio · lámina 24}{¿Cuál es el control negativo de esa cadena?}
\A{Repetir la predicción con $c_s$ en lugar de $v_{f\perp}$: el resultado
colapsa a $\hat\omega\approx0.016$, seis veces por debajo de lo medido.
Eso demuestra que el acuerdo del 8\% no es genérico ---no cualquier
velocidad lo produce--- y que es la presión magnética del campo guía la
que sostiene la inestabilidad en nuestro régimen: sin ella el flujo sería
demasiado supersónico respecto de la restitución puramente acústica.}

\Q{nivel medio · lámina 24}{¿Dónde queda su flujo respecto del umbral de
estabilización relativista?}
\A{El criterio (Königl 1980; Bodo 2004; Chow 2023) es que la discontinuidad
es inestable si el Mach relativista construido con la velocidad de
restitución cumple $0<M_{re}<\sqrt2$. El nuestro es $M_{re}=0.60$:
holgadamente inestable, lejos del umbral $\sqrt2\approx1.41$. Importa
porque la extrapolación de la asíntota no roza ninguna frontera de
estabilidad: el plateau ideal es genuino, no un borde de supresión.}

\Q{nivel alto · lámina 24}{¿Por qué no resolvieron la relación de dispersión
RMHD completa?}
\A{La dispersión compresible magnetizada relativista para el perfil suave
es un polinomio de grado 8 en la frecuencia con coeficientes dependientes
del perfil (Osmanov 2008; Chow 2023): resolverla con nuestros parámetros es
un proyecto en sí mismo, declarado como trabajo futuro. Para el contraste
de este trabajo empleamos la cadena validada Michalke $\to$ sustitución
escalar, y por eso presentamos el resultado como \emph{consistencia
cuantitativa fuerte} ---8\%, sin parámetros libres--- y no como
validación cerrada. La asíntota, además, es extrapolación: el dato máximo
alcanza el 85\% del plateau.}

%========================= PARTE C: categorías 11–15 =========================

\section{Régimen no lineal y variables globales}

\Q{nivel básico · lámina 25}{Enuncie su resultado no lineal principal.}
\A{El máximo de la enstrofía de perturbación escala como ley de potencia
con la conductividad: $\Omega^{\max}_{zp}\propto\sigma^{1.22}$ con
$R^2=0.997$ sobre las 29 simulaciones con pico identificable. El punto de
comparación es el piso teórico de una lámina de corriente única de
Sweet--Parker, que daría exponente $1/2$: nuestro exponente lo excede en un
factor $\sim2.4$.}

\Q{nivel medio · lámina 25}{¿Qué es el piso de Sweet--Parker y por qué es la
referencia?}
\A{En el modelo de Sweet--Parker (Parker 1957) una lámina de corriente
laminar de longitud $L$ tiene espesor $\delta\sim L\,S^{-1/2}$: la
vorticidad concentrable en ella escala entonces como $\sigma^{1/2}$. Es la
referencia natural porque es el régimen de reconexión resistiva más
sencillo y porque el escalamiento del espesor es robusto frente a
correcciones relativistas (Lyubarsky 2005): si el sistema tuviera una sola
lámina laminar, mediríamos $1/2$.}

\Q{nivel medio · lámina 25}{¿Cómo interpretan el exceso del exponente y con
qué grado de certeza?}
\A{Como hipótesis fenomenológica, declarada como tal: el exceso sugiere que
la lámina primaria se fragmenta por el modo de desgarro (\emph{tearing},
Furth, Killeen \& Rosenbluth 1963), generando plasmoides que multiplican la
superficie disipativa efectiva; la firma es consistente con lo reportado
por Nastro et al.\ (2022). No lo presentamos como constante derivada:
demostrarlo exige resolver la jerarquía de sub-láminas, que a nuestras
resoluciones está en el borde de la malla ---de ahí la propuesta de
esquemas de cuarto orden (Mignone 2024) como trabajo futuro.}

\Q{nivel básico · lámina 26}{Resuma la firma electromagnética de $\sigma$ en
las variables globales.}
\A{Tres tendencias monótonas. El trabajo electromagnético disipado
acumulado crece de $0.03$ a $0.44$ a lo largo del barrido: más
conductividad permite más violencia hidrodinámica que procesa más energía
magnética. La amplificación de la energía magnética propia decae del
15\% a 0 hacia el ideal: el congelamiento impide reorganizar la
topología. Y las estructuras fuera del plano ($f_{Vz}$) tienen su máximo
$0.37$ en $\sigma\approx1400$, dentro de la transición.}

\Q{nivel alto · lámina 26}{¿No es contradictorio que la disipación acumulada
\emph{crezca} con $\sigma$, si $\eta=1/\sigma$ decrece?}
\A{Es la paradoja aparente que más nos gusta explicar. La potencia óhmica
local es $\sim\eta J^2$: al subir $\sigma$, $\eta$ cae pero las láminas de
corriente se adelgazan y $J$ crece más que compensando ($J\propto
\delta^{-1}$ con $\delta\sim S^{-1/2}$). Además la inestabilidad es más
vigorosa ---$\gamma$ y $\Omega^{\max}$ crecen--- y procesa más energía
magnética por unidad de tiempo. El integrando pequeño sobre estructuras
intensas y duraderas gana: la disipación acumulada crece con $\sigma$ en
todo nuestro rango.}

\Q{nivel medio · lámina 26}{¿Por qué el plateau de $J_{\max}$ es el
argumento decisivo contra la disipación numérica?}
\A{Porque $J_{\max}$ mide el gradiente más agudo del campo: si lo limitara
la malla, crecería sin saturar al aumentar $\sigma$ hasta el límite de
resolución en todos los casos. Lo observado es una saturación ---plateau---
a partir de la conductividad en que la capa resistiva física es más ancha
que el error de truncamiento: desde ahí la difusividad \emph{física}
gobierna el espesor. Que la zona de transición completa esté dentro del
plateau acredita que sus resultados no están contaminados por disipación
artificial.}

\section{Campaña magnética (objetivo 4)}

\Q{nivel básico · lámina 27}{Campaña A: ¿qué suprime la inestabilidad al
aumentar $B_0$?}
\A{La presión magnética, no la tensión. Al pasar $B_0$ de $1.0$ a $2.0$
(a $\sigma=10^4$), $\gamma$ cae de $0.81$ a $0.38$ y $\Omega^{\max}$ de
$87$ a $25$; pero el flujo sigue siendo súper-alfvénico paralelo
($M_{A,\parallel}$ de $9.4$ a $6.4$), de modo que la tensión es incapaz de
frenar la cizalla. El agente es la magnetización creciente: $\beta$ cae de
$0.99$ a $0.25$, el medio se rigidiza por presión magnética y el
reservorio de energía libre se procesa menos eficientemente.}

\Q{nivel medio · lámina 27}{Explique la ``paradoja'' de $\theta=0^\circ$ en
la campaña B.}
\A{Con campo guía puro no hay tensión sobre el modo, y sin embargo la
enstrofía es mínima: parecería que quitar el freno diera menos vorticidad.
La resolución es de Mizuno (2013): en 2D el vórtice solo puede estirar
líneas de campo que tengan componente en el plano; sin $B_y$ no hay
estiramiento, luego no hay dínamo local ni reconexión que realimenten la
vorticidad. La pequeña $B_y$ de los demás ángulos actúa como semilla de
amplificación. La paradoja se disuelve al distinguir el papel de freno
(tensión sobre el modo) del papel de alimentación (estiramiento en el
plano).}

\Q{nivel medio · lámina 27}{En $\theta=90^\circ$ el máximo local de
enstrofía crece pero la integrada decrece. ¿Contradicción?}
\A{No: es un efecto del observable. Con tensión máxima ($B\parallel
\mathbf k$) el crecimiento global está suprimido y la enstrofía
\emph{integrada} decrece, como manda la teoría. Pero la misma tensión
organiza capas de reconexión finas e intensas cuyo máximo \emph{local} es
alto. Un observable global y uno local miden física distinta; reportarlos
juntos desambigua la aparente paradoja y es una lección metodológica del
trabajo.}

\Q{nivel básico · lámina 29}{¿Qué papel juega la campaña C en el diseño?}
\A{Es el experimento de control. Al rotar el campo en el plano $x$--$z$ se
mantiene $\mathbf k\cdot\mathbf B=0$ para todo ángulo: no hay tensión sobre
el modo, y además $|\mathbf B|$ ---y con él la presión magnética--- es
constante con $\theta$. Cualquier variación observada es entonces
puramente geométrica (reorientación de $B_x$, topología de reconexión).
Comparar B (tensión creciente) contra C (tensión nula) aísla el efecto de
la tensión con la presión controlada.}

\Q{nivel medio · lámina 29}{¿Qué encontraron sin tensión y qué concluyen?}
\A{Sin tensión no hay fase lineal limpia: el crecimiento pierde el carácter
exponencial y la mezcla se vuelve difusa, de forma idéntica a
$\sigma=6000$ y $10\,000$. Con tensión (campaña B) la fase exponencial
existe. Conclusión: la tensión magnética es el estabilizador ---y
organizador--- primario de la KHI en este régimen; la presión modula la
amplitud, pero es la tensión la que estructura la dinámica. Además, sin el
intercambio elástico guiado por tensión, la anticorrelación
cinético-magnética se intensifica hasta $r\approx-0.98$: conversión casi
uno a uno.}

\Q{nivel básico · lámina 28}{Describa la anti-fase energética y qué decide
la resistividad.}
\A{Las partes oscilatorias de $E_{\rm kin}$ y $E_{\rm mag}$ oscilan en
contrafase ciclo a ciclo (correlaciones de $-0.6$ a $-0.9$ en B; hasta
$-0.98$ en C): el vórtice estira campo ---dínamo, cinética a magnética--- y
la tensión devuelve ---rebote elástico---. La resistividad decide el
destino: a $\sigma=10^4$ el intercambio es cuasi-reversible; a
$\sigma=6000$ la disipación óhmica desvía parte de cada ciclo a calor,
hasta un orden de magnitud más a $\theta=90^\circ$, donde las capas de
reconexión son más intensas.}

\Q{nivel medio · lámina 28}{¿Por qué usar el balance energético como
observable en la campaña magnética?}
\A{Porque la geometría del campo enmascara la fase lineal: con tensión
fuerte o sin fase exponencial no hay pendiente que ajustar, y $\gamma$
deja de ser un observable útil. Las energías, en cambio, existen siempre,
son integrales (robustas) y su intercambio y disipación son exactamente la
física que el objetivo 4 pregunta. La lección metodológica: el observable
debe elegirse según el régimen, no imponerse por costumbre.}

\Q{nivel medio · lámina 28}{¿Qué control de calidad interno acompaña estos
resultados?}
\A{La conservación de la energía total: en todas las simulaciones de la
campaña el error relativo de $E_{\rm tot}$ es $\lesssim10^{-3}$ en el
horizonte completo. Es una validación independiente del esquema ---la
conservación no se impone a las energías parciales, emerge de la
formulación conservativa--- y garantiza que la anti-fase no es un artefacto
de pérdida numérica de energía.}

\section{Estadística y validación del ajuste}

\Q{nivel básico · respaldo B8}{¿Qué es el AIC y qué significa su
$\Delta\mathrm{AIC}=-126$?}
\A{El criterio de información de Akaike, $\mathrm{AIC}=2k-2\ln\hat L$,
penaliza la verosimilitud con el número de parámetros: castiga el
sobreajuste. Comparando modelos anidados sobre los \emph{mismos} datos, el
sigmoide con desplazamiento mejora el AIC en 126 unidades respecto del de
tres parámetros; en la escala de Burnham \& Anderson (2002), diferencias
mayores que 10 ya son evidencia decisiva. El parámetro extra paga su costo
con creces.}

\Q{nivel medio · respaldo B8}{Describa su validación fuera de muestra.}
\A{Los 9 puntos más resistivos ($\sigma<1600$) no participan del ajuste
---allí no hay fase lineal limpia y su peso sería dudoso---. Usados como
conjunto de prueba, el sigmoide con desplazamiento los predice con RMSE
$0.025$ contra $0.039$ del modelo sin desplazamiento: 36\% mejor. Es la
prueba más exigente disponible: capacidad predictiva en el régimen donde el
modelo no fue entrenado.}

\Q{nivel alto · respaldo B8}{Sus parámetros $C$ y $\gamma_0$ están
degenerados. ¿Cómo lo manejan?}
\A{La correlación es $r(C,\gamma_0)=-0.997$: el ajuste puede intercambiar
desplazamiento por amplitud casi libremente, así que ninguno de los dos es
confiable por separado. La combinación que la degeneración deja invariante
es la suma $C+\gamma_0$ ---la asíntota---, que varía menos del 4\% entre
variantes del ajuste. Por eso el resultado reportado es la asíntota, con
su incertidumbre, y no los parámetros individuales.}

\Q{nivel medio · respaldo B8}{¿Qué revelaron los residuos del ajuste y qué
hicieron al respecto?}
\A{Autocorrelación: rachas de residuos del mismo signo, confirmadas con un
test de rachas. Significa que el sigmoide no captura toda la estructura
fina de $\gamma(\sigma)$ y que los errores formales de la covarianza
\emph{subestiman} la incertidumbre real. La respuesta metodológica fue no
confiar en el error formal: $\sigma_{\rm crit}$ se estimó con cuatro
métodos independientes y se reporta la banda completa $[1400,7000]$.}

\Q{nivel medio · lámina 22}{¿Cómo asignaron errores a los cuatro
estimadores?}
\A{Para cada observable, el estimador es el punto de inflexión respecto de
$\log_{10}\sigma$; el error es el máximo entre dos medidas conservadoras:
la dispersión jackknife ---repetir la estimación quitando un punto a la
vez--- y el semiancho de la región de inflexión. Así, estimadores basados
en pocos puntos o con inflexión ancha ($\Omega^{\max}$: $\pm2175$) llevan
honestamente errores grandes, y el promedio inter-método no queda dominado
por una precisión ilusoria.}

\Q{nivel alto · lámina 21}{¿Consideraron modelos alternativos al sigmoide?}
\A{Sí, dentro de la familia razonable para una transición suave entre dos
plateaus en $\log\sigma$: el sigmoide simple (descartado por AIC y
hold-out), y leyes de potencia con saturación, que carecen de plateau
resistivo y describen mal el extremo bajo. El sigmoide con desplazamiento
es el modelo mínimo que captura ambos regímenes con parámetros
interpretables ---asíntota, centro y anchura de la transición---. No
afirmamos que sea la forma funcional verdadera: es el descriptor
fenomenológico más parsimonioso compatible con los datos.}

\section{Epistemología, limitaciones y trabajo futuro}

\Q{nivel alto · respaldo B9}{Desde Liouville hasta su código: justifique la
cadena de aproximaciones.}
\A{El teorema de Liouville conserva la densidad de probabilidad en el
espacio de fase; para un plasma, la jerarquía BBGKY la reduce a la
ecuación cinética de una partícula: Vlasov--Boltzmann con integrales de
colisión. Tomando momentos (densidad, momento, energía) y cerrando con
equilibrio termodinámico local ---válido si el camino libre medio es corto
frente a la escala hidrodinámica--- se llega al fluido único con
coeficientes de transporte, entre ellos $\eta$. Nuestro código integra ese
nivel. La cadena es explícita y cada eslabón tiene su condición de validez
declarada; el respaldo B9 la recorre paso a paso.}

\Q{nivel alto · lámina 31}{¿Cuándo colapsa epistemológicamente su
descripción de fluido único?}
\A{Cuando la escala disipativa relevante deja de ser fluida: si la capa
resistiva $\delta\sim S^{-1/2}L$ se hace comparable al girorradio iónico o
al camino libre medio, el cierre colisional pierde sentido y harían falta
efecto Hall, presiones anisótropas o descripción cinética
(Particle-in-Cell). En nuestro régimen declarado ---plasma colisional
denso, $\nu_{\rm col}\gg\omega_{\rm cicl}$--- la jerarquía está del lado
seguro. Lo importante metodológicamente es que el límite está identificado
\emph{antes} de interpretar los resultados, no después.}

\Q{nivel medio · lámina 31}{Enumere las limitaciones declaradas del
trabajo.}
\A{Cinco. Bidimensionalidad: sin modos 3D ni cascada turbulenta plena. EoS
politrópica con $\Gamma=4/3$ fijo: sin composición ni radiación.
Conductividad escalar constante: sin Hall, anisotropía ni dependencia
térmica. Resolución: a las $\sigma$ más altas las láminas rozan la malla,
lo que condiciona el exponente $1.22$. Y la asíntota ideal es
extrapolación: el dato máximo alcanza el 85\% del plateau. Cada una está
declarada junto al resultado que condiciona.}

\Q{nivel medio · lámina 31}{¿Cuál es su programa de trabajo futuro y su
prioridad?}
\A{En orden: (1) medir el plateau directamente con 2--3 simulaciones a
$\sigma\sim2$--$3\times10^4$ y mayor resolución, convirtiendo la
extrapolación en medición; (2) resolver la dispersión RMHD compresible
completa de grado 8 en nuestros parámetros, cerrando el contraste teórico;
(3) mapa fino de la frontera tensión-presión con más ángulos y
conductividades; (4) extensión 3D con esquemas de cuarto orden (Mignone
2024) para resolver la jerarquía de sub-láminas y poner a prueba la
hipótesis de \emph{tearing}; y a largo plazo, híbridos fluido-cinéticos
para la frontera no colisional.}

\Q{nivel alto}{¿Qué resultado numérico falsaría su modelo sigmoide?}
\A{Varios, y son medibles: si al simular $\sigma\sim2$--$3\times10^4$ la
tasa siguiera creciendo en lugar de saturar en $1.03\pm0.05$, la asíntota
---y con ella el contraste con la teoría ideal--- quedaría refutada; si al
duplicar la resolución el exponente no lineal se moviera sustancialmente
de $1.22$, la firma de \emph{tearing} sería un artefacto; y si el plateau
de $J_{\max}$ se desplazara con la malla, la separación disipación
física/numérica caería. El diseño deja definidos los experimentos que
podrían derribar cada afirmación: esa es nuestra noción operativa de
falsabilidad.}

\Q{nivel alto}{¿No es la simulación una ``caja negra'' entre la teoría y
usted?}
\A{No en este trabajo, por tres razones. Conocemos el interior: qué
polinomio resuelve la recuperación de primitivas, qué ecuación obedece el
error de divergencia, qué límite asintótico preserva el integrador.
Controlamos los extremos: el sistema recupera analíticamente RMHD ideal,
MHD clásica y Maxwell, y el solucionador lineal reproduce a Michalke. Y
medimos la frontera entre lo físico y lo numérico (plateau de $J_{\max}$).
La simulación funciona aquí como un experimento controlado sobre unas
ecuaciones que no admiten solución cerrada: un puente entre la teoría de
campos y la fenomenología, no una caja negra.}

\Q{nivel medio}{¿Por qué su trabajo usa ``experimento numérico'' como
categoría? ¿No es un oxímoron?}
\A{Lo usamos con precisión: hay variable independiente ($\sigma$, y luego
$B_0,\theta$), variables dependientes definidas antes de medir, protocolo
fijo (misma malla, misma ventana, mismos criterios), controles (límites
analíticos, conservación de $E_{\rm tot}$, validación de Michalke) y
condiciones de exclusión declaradas. Que el aparato sea un sistema de
ecuaciones discretizado y no un túnel de viento no cambia la lógica
inferencial; cambia el tipo de error sistemático, que por eso dedicamos
tanto espacio a acotar.}

\Q{nivel medio · lámina 30}{¿Por qué sus conclusiones no afirman
``objetivo cumplido''?}
\A{Por respeto al proceso de evaluación: las conclusiones enuncian los
resultados en correspondencia uno a uno con los objetivos ---qué se midió,
con qué incertidumbre y con qué alcance--- y la valoración del cumplimiento
corresponde al jurado y al director. Es además coherencia epistemológica:
un trabajo que insiste en declarar límites no puede autocertificarse.}

\section{Sobre la bibliografía citada}

\Q{nivel medio}{¿Qué toman de Komissarov (2007) y de Palenzuela et al.\
(2009)?}
\A{Son los dos pilares de la formulación. De Komissarov: el esquema RRMHD
multidimensional con el sistema aumentado ---telégrafo para las
ligaduras--- y el diagnóstico de que la reconexión ideal-numérica es
incontrolada. De Palenzuela: el tratamiento IMEX del término rígido, la
filosofía ``beyond ideal MHD'' y el argumento del suavizado parabólico de
la discontinuidad tangencial, que fundamenta por qué $\gamma(\sigma)$ debe
medirse numéricamente.}

\Q{nivel básico}{¿Cuál es el linaje del código Cueva en la literatura?}
\A{Miranda-Aranguren, Aloy \& Aloy-Torás (2014, IAU) presentan la
construcción del código RRMHD no ideal; Miranda-Aranguren et al.\ (2018)
publican el solucionador HLLC para RRMHD y las pruebas de convergencia y
choques del esquema, incluyendo la observación ---central para nosotros---
de que MP5 compensa la difusión de contacto de HLL. Nuestro trabajo hereda
esa infraestructura validada y la aplica al problema físico del barrido en
conductividad.}

\Q{nivel medio}{¿Qué papel juegan Dedner (2002), Munz (2000) y Lee (2004) en
su capítulo de divergencias?}
\A{Munz introduce la limpieza hiperbólica generalizada para Maxwell
---nuestro sector electromagnético---; Dedner la formula para MHD con la
pareja hiperbólico-parabólica que usamos en la práctica; y Lee aporta la
estructura lagrangiana y de simetrías de la limpieza, que es la base del
respaldo B3: el GLM no es un truco, tiene origen variacional y por eso
respeta causalidad y disipación de forma controlada.}

\Q{nivel medio}{¿Por qué citan a Rueda-Ramírez et al.\ (2023), que es de
métodos entropía-estables?}
\A{Porque justifica los términos de Godunov--Powell de nuestro vector
fuente: ese trabajo demuestra que, en presencia de errores de divergencia
transitorios, la variante con términos no conservativos proporcionales a
las ligaduras es compatible con una desigualdad de entropía. Es la
referencia moderna que convierte una corrección históricamente heurística
en un requisito termodinámico riguroso.}

\Q{nivel básico}{¿Qué aportan Suresh \& Huynh (1997), Toro (2009) y Harten,
Lax \& van Leer (1983)?}
\A{La caja de herramientas hiperbólica. Harten--Lax--van Leer: el marco de
solucionadores aproximados de dos ondas del que sale HLL. Toro: el
tratado de referencia sobre solucionadores de Riemann y las condiciones de
Rankine--Hugoniot en volúmenes finitos. Suresh \& Huynh: la reconstrucción
MP5 con preservación de monotonicidad, clave para retener extremos de
vorticidad sin oscilaciones de Gibbs.}

\Q{nivel medio}{¿Y Pareschi \& Russo (2005) y Aloy \& Cordero-Carrión
(2016)?}
\A{Pareschi \& Russo definen los esquemas IMEX Runge--Kutta fuertemente
estables (SSP) para sistemas hiperbólicos con relajación rígida: nuestro
integrador pertenece a esa familia y hereda sus propiedades de estabilidad
y preservación asintótica. Aloy \& Cordero-Carrión desarrollan la
alternativa mínimamente implícita (MIRK) para RRMHD: la citamos como la
otra rama contemporánea del mismo problema, contra la cual se entiende la
elección IMEX de Cueva.}

\Q{nivel básico}{¿Por qué Mizuno (2013) es el origen de su montaje?}
\A{Define el problema de referencia: estudia el rol de la ecuación de
estado en la KHI resistiva relativista con la doble capa de cizalla que
adoptamos (nuestro ``tipo Mizuno''), incluyendo el argumento de $B_y$ como
semilla de amplificación que resuelve la paradoja de $\theta=0^\circ$.
Partir de un montaje publicado hace comparables nuestros resultados y
separa nuestra contribución ---el barrido sistemático en $\sigma$ y la
campaña de orientación--- de la definición del escenario.}

\Q{nivel medio}{Sitúe a Michalke (1964), Landau (1944), Königl (1980), Bodo
et al.\ (2004), Osmanov et al.\ (2008) y Chow et al.\ (2023) en su cadena
teórica.}
\A{Michalke: la solución clásica del perfil $\tanh$ inviscido, nuestro
techo y validación del solucionador. Landau: la estabilización compresible
del \emph{vortex sheet} (umbral $\sqrt2$), el origen histórico del
criterio. Königl: su extensión relativista. Bodo: el criterio y la
evolución 3D en cizalla magnetizada relativista. Osmanov: la dispersión
RMHD plana completa (grado 8), la referencia de lo que dejamos como
futuro. Chow: la teoría lineal relativista magnetizada reciente, de donde
tomamos ``presión pero no tensión'' y la sustitución $c_s\to v_{f\perp}$.}

\Q{nivel medio}{¿Qué aportan Parker (1957), Furth et al.\ (1963), Lyubarsky
(2005), Nastro et al.\ (2022), Lecoanet et al.\ (2016) y Mignone et al.\
(2024) al análisis de resultados?}
\A{Parker: el modelo Sweet--Parker, nuestro piso $\sigma^{1/2}$. Furth,
Killeen \& Rosenbluth: el modo \emph{tearing}, la hipótesis del exceso.
Lyubarsky: la robustez relativista del escalamiento del espesor. Nastro:
el estudio KHI resistivo con el que comparamos metodología (enstrofía,
factor 2) y la firma de fragmentación. Lecoanet: la no convergencia de la
KHI 2D ideal, pilar del argumento de buen planteamiento. Mignone: los
esquemas de cuarto orden que harían resoluble la jerarquía de sub-láminas.}

\Q{nivel básico}{¿Qué papel cumplen los textos de fondo (Misner--Thorne--%
Wheeler, Schutz, Rezzolla--Zanotti, Goedbloed, Nakahara, Baez--Muniain,
Bachman, Hehl--Obukhov, Burnham--Anderson)?}
\A{Fijan lenguaje y convenciones. Misner y Schutz: cálculo tensorial y
signatura. Rezzolla--Zanotti: la formulación de Valencia y la
hidrodinámica relativista numérica. Goedbloed: la MHD clásica de
referencia. Nakahara, Baez--Muniain, Bachman y Hehl--Obukhov: el lenguaje
de formas diferenciales con el que presentamos $F=dA$, la nilpotencia y la
electrodinámica geométrica (respaldos B1--B3). Burnham--Anderson: la
metodología de selección de modelos por AIC. Citarlos delimita qué es
estándar y qué es contribución propia.}

\Q{nivel alto}{Si el jurado le pide una referencia que conecte su trabajo
con la física que él enseña (electrodinámica y mecánica estadística),
¿cuál elige y cómo la presenta?}
\A{Dos. Hehl \& Obukhov para electrodinámica: presenta las ecuaciones de
Maxwell como enunciados geométricos sobre formas diferenciales
---conservación de flujo y de carga--- independientes de la métrica, el
puente exacto entre la electrodinámica de pregrado y nuestro $F=dA$ con
Bianchi. Y para mecánica estadística, la cadena
Liouville$\to$Vlasov$\to$momentos del respaldo B9, con Goedbloed como
texto de referencia del cierre fluido: es literalmente el contenido del
teorema de Liouville aplicado hasta donde nace nuestra $\eta$.}


\end{document}
